% !TEX root = ../main.tex
% =====================================================================
\section{Zeiger}
% =====================================================================
\begin{symtab}
$U$ & Spannung & [V] \\
$I$ & Strom & [A] \\
$\hat{I}$ & Stromamplitude & [A] \\
$L_{ph}$ & Phaseninduktivität & [H] \\
$\psi$ & Flussverkettung & [Wb] \\
$\omega$ & Kreisfrequenz & [rad/s] \\
$f$ & Frequenz & [Hz] \\
$t$ & Zeit & [s] \\
$\varphi_i$ & Anfangsphase des Stroms & [rad] \\
$\theta$ & Raumwinkel (Wicklungsachse) & [rad] \\
\end{symtab}

\subsection{Drehfeldzeiger}
\fig[0.6]{drehfeldzeiger}

Drei räumlich um $120^\circ$ versetzte Spulen mit zeitlich um $-120^\circ$
phasenverschobenen Strömen erzeugen ein Drehfeld konstanter Amplitude.\\
Für die drei um $120^\circ$ phasenverschobenen Ströme gilt
\[
i_k(t) = \hat{I}\cos(\omega t + \varphi_k), \qquad \psi_k(t) = L_{ph} i_k(t), \quad k \in \{a,b,c\}.
\]
Die Flussverkettungen der drei Phasen bilden durch geometrische Addition
ein Drehfeld. Die reelle Achse des Koordinatensystems fällt mit der
Wicklungsachse von Phase a zusammen. Die Addition der Momentanwerte liefert
den resultierenden Flussverkettungs-Drehzeiger.

\columnbreak
\fig[0.6]{drehfeldzeiger_fluss}

Mit dem Konventionsfaktor $2/3$ ergibt sich für den Flussverkettungs-Drehzeiger
\[
\vec{\psi} = \frac{2}{3} \sum_{k=a,b,c} \psi_k e^{j\theta_k} = \hat{I} L_{ph} e^{j(\omega t + \varphi_i)}
\]

\subsubsection{Drehfeldzeigertransformation: Dreiphasenebene $\Rightarrow$ Drehzeiger}
Dreiphasige Grössen lassen sich analog zur Flussverkettung als Drehzeiger
in der komplexen Ebene darstellen. Die Transformation gilt allgemein für
Flussverkettungen, Ströme, Spannungen sowie beliebige Dreiphasengrößen $x$.
\begin{gather*}
\vec{\psi}(t) = \frac{2}{3}\Big[\psi_a(t) + \psi_b(t) e^{j120^\circ} + \psi_c(t) e^{-j120^\circ}\Big] \\
\vec{i}(t) = \frac{2}{3}\Big[i_a(t) + i_b(t) e^{j120^\circ} + i_c(t) e^{-j120^\circ}\Big] \\
\vec{u}(t) = \frac{2}{3}\Big[u_a(t) + u_b(t) e^{j120^\circ} + u_c(t) e^{-j120^\circ}\Big] \\
\vec{x}(t) = \frac{2}{3}\Big[x_a(t) + x_b(t) e^{j120^\circ} + x_c(t) e^{-j120^\circ}\Big]
\end{gather*}
Die Drehzeigerdarstellung ermöglicht eine vereinfachte Analyse gegenüber
der Beschreibung mit einzelnen Phasengrössen. Für symmetrische,
sinusförmige Dreiphasensysteme besitzt der Drehzeiger eine konstante Länge
und rotiert mit der Frequenz der Phasengrössen. Seine Spitze beschreibt
dabei einen Kreis in der komplexen Ebene.

\subsubsection{Drehfeldzeigertransformation: Drehzeiger $\Rightarrow$ Zeiger}
Durch Transformation in ein mit $\omega_K$ rotierendes Koordinatensystem wird der
Drehzeiger zu einem zeitlich konstanten Zeiger.
\begin{gather*}
X(t) = \vec{x}(t)\, e^{-j\int \omega_K\, dt}, \\
X(t) = \vec{x}(t)\, e^{-j\omega_K t} \qquad (\omega_K = \text{konstant}), \\
X(t) = X_x + j\,X_y \;(= X_d + j\,X_q).
\end{gather*}
\begin{achtung}
$\omega$ beschreibt die physikalische Drehung\\
$\omega_K$ die Drehung des Bezugssystems.\\
Gleich sind sie \textbf{nur} im synchronen Koordinatensystem
\end{achtung}

\columnbreak
\subsubsection{Drehfeldzeigertransformation: Drehzeiger $\Rightarrow$ Phasengrösse}
Die Rücktransformation ist eindeutig nur bei Berücksichtigung der gleichphasigen Komponente $x_0$.
\vspace{2em}
\begin{align*}
x_a &= x_\alpha + x_0, \\
x_b &= \tfrac{1}{2}(\sqrt{3}x_\beta - x_\alpha) + x_0 \\
x_c &= \tfrac{1}{2}(-\sqrt{3}x_\beta - x_\alpha) + x_0
\end{align*}
\vspace{1em}
